% titel: Blatt 3, A1: Kinetische und potentielle Energie
% skript: Kap. Energieerhaltung

% Energieerhaltung, Federenergie, Uebergang kinetisch-potentiell.

Ohne Reibung bleibt die mechanische Energie $E=E_{\rm kin}+E_{\rm pot}$
konstant; an einer Feder kommt $E_{\rm pot}=\tfrac12 kx^2$ hinzu.

\begin{werkzeugkasten}[Werkzeuge]
\textbf{Mechanische Energieerhaltung} (reibungslos):
\begin{equation}
  E=\tfrac12 m v^2+mg\,h=\text{const.},
\end{equation}
wobei $h$ gegen eine frei waehlbare Nulllage gemessen wird. Fuer zwei
Punkte 1,2:
\begin{equation}
  \tfrac12 m v_1^2+mgh_1=\tfrac12 m v_2^2+mgh_2\;.
\end{equation}

\textbf{Federenergie} (Auslenkung $x$ der Feder):
\begin{equation}
  E_{\rm Fed}=\tfrac12 k x^2.
\end{equation}
Arbeit der Kraft an der Feder: $W=\tfrac12kx^2$.

\textbf{Potentielle Energie $E_{\rm pot}(x)$} ist die Ortsfunktion jeder
konservativen Kraft. Man skizziert sie stückweise: Gravitation
$mgh(x)$ plus Feder $\tfrac12kx^2$ (nur dort, wo die Feder greift).
\end{werkzeugkasten}

\begin{werkzeugkasten}[Vorgehen]
\begin{itemize}
  \item Zwischen je zwei Bahnpunkten Energieerhaltung ansetzen und nach
        $v$ aufloesen.
  \item Zum tiefsten Punkt (B) alles kinetische Vorwaerts-Hohengefälle
        nehmen.
  \item Federeindruck: kinetische Energie beim Treffen der Feder in
        $\tfrac12kx^2$ umwandeln.
  \item Fuer die Rueckfahrt die neue Nulllage und beide Energieformen
        zusammenfuehren.
\end{itemize}
\end{werkzeugkasten}

\begin{achtung}[Typische Fallen]
\begin{itemize}
  \item Die potentielle Energie haengt nur von der \emph{relative hohe}
        (Hohendifferenz) ab, nicht vom Weg.
  \item Nullpunkt der Federauslenkung (im entspannten Zustand) und der
        Horizontalen-Koordinate nicht verwechseln.
  \item Beim Rueckweg gilt dieselbe Energiebilanz; $E_{\rm pot}(x)$ waechst
        beim Federn um $\tfrac12kx^2$.
\end{itemize}
\end{achtung}
